\subsection{2.2 连续运行能力}

连续运行能力占 20\%，在北京或天津规定的高速/快速路开放道路上跑，更像真实用户体验测试。

下面有三个指标：

\textbf{通行成功率}包括：变道成功率、汇入匝道成功率、汇入主路成功率。比如变道成功要求系统自动控制转向、提前打转向灯、按交通法规合理完成变道。未打灯、变道中途返回、打灯后不变道，都算失败。

\textbf{系统提醒频次}看的是每 100 km 里，系统提醒驾驶员介入、功能降级、功能退出的次数。高速路要求比快速路更严格，例如白天/夜晚高速路系统提醒频次 ≤ 15 次/100 km 才能拿 100 分。

\textbf{主动介入频次}看的是驾驶员因为紧急情况主动接管的次数。高速路同样更严格，白天/夜晚高速路主动介入频次 ≤ 8 次/100 km 才能拿 100 分。

\textbf{舒适度}分横向和纵向。横向看 2-3 m/s²、>3 m/s² 横向加速度出现频次；纵向看 2-4 m/s²、>4 m/s² 纵向减速度出现频次。出现越频繁，说明系统越容易急打方向或急刹，得分越低。

\subsubsection{2.2.1 通行成功率}

通行成功率看系统自己能不能完成路线任务，包含变道成功率、汇入匝道成功率和汇入主路成功率，权重分别为 40\%、40\%、20\%。

计算方式为：成功率 = 成功次数 / （成功次数 + 失败次数） × 100\%。

变道成功要求系统提前激活转向灯、合理控制车辆，并按照交通法规完成变道。未点亮转向灯、变道中途返回原车道、点亮转向灯后未变道等情况，都计为变道失败。

通行成功率评分阈值：

\subsubsection{2.2.2 系统提醒/主动介入频次}

系统提醒/主动介入频次看系统在长距离开放道路运行时有多稳定。系统提醒频次统计每 100 km 内系统提醒驾驶员介入、功能降级、功能退出的次数；主动介入频次统计每 100 km 内驾驶员因紧急情况主动接管的次数。

计算方式分别为：系统提醒频次 = 100 × 提醒次数 / 测试里程；主动介入频次 = 100 × 主动介入次数 / 测试里程。

其中，道路施工、交通事故，以及变道失败、汇入匝道/主路失败导致的提醒或接管，不计入本项频次，因为这些已经在对应指标中评价。

系统提醒/主动介入频次评分阈值：

\subsubsection{2.2.3 车辆行驶舒适度}

车辆行驶舒适度看系统有没有频繁急打方向或急刹，分为横向控制舒适度和纵向控制舒适度，权重各 50\%。

横向控制舒适度统计每 100 km 内横向加速度事件：2 m/s² ＜ a ≤ 3 m/s² 计一次，a ＞ 3 m/s² 按两倍计入。纵向控制舒适度统计每 100 km 内纵向减速度事件：2 m/s² ＜ a ≤ 4 m/s² 计一次，a ＞ 4 m/s² 按两倍计入。

车辆行驶舒适度评分阈值：
